% Work-MAT, hodina 56. Navazuje na 52--55; zdroj ucebnice Algebraické výrazy,
% oddily Lomene vyrazy a Operace s mnohocl. Metodicka inspirace Realisticky.cz 010905.
\documentclass[11pt,a4paper]{article}
% Spolecny styl pro ucitelske pripravy MAT (Work/MAT)
% Verze 1.1, 2026-10-23
\NeedsTeXFormat{LaTeX2e}
\ProvidesPackage{priprava}[2026/10/23 v1.1 Sablona ucitelskych priprav MAT]

% Zakladni typografie a cestina
\RequirePackage[T1]{fontenc}
\RequirePackage[utf8]{inputenc}
\RequirePackage[czech]{babel}
\RequirePackage{lmodern}
\RequirePackage{microtype}
\RequirePackage{amsmath,amssymb}
\RequirePackage{xcolor}
\RequirePackage{enumitem}
\RequirePackage[most]{tcolorbox}
\RequirePackage{tabularx,booktabs,array}
\RequirePackage{fancyhdr}
\RequirePackage{needspace}
\RequirePackage{geometry}
\RequirePackage[hidelinks]{hyperref}

\geometry{a4paper,top=17mm,bottom=18mm,left=23mm,right=23mm}
\definecolor{notegray}{gray}{0.39}
\setlength{\parindent}{0pt}
\setlength{\parskip}{2.5pt plus 1pt minus 0.5pt}
\setlist[itemize]{topsep=2pt,itemsep=1pt,parsep=0pt,leftmargin=1.45em}
\setlist[enumerate]{topsep=2pt,itemsep=2pt,parsep=0pt,leftmargin=1.65em}
\setcounter{secnumdepth}{0}
\renewcommand{\arraystretch}{1.12}

% Zahlavi: vlevo rocnik | globalni tema, vpravo cislo: nazev hodiny
% Zapati zustava zcela prazdne.
\pagestyle{fancy}
\fancyhf{}
\renewcommand{\headrulewidth}{0.35pt}
\renewcommand{\footrulewidth}{0pt}
\setlength{\headheight}{14pt}
\newcommand{\nastavPripravu}[4]{%
  \fancyhead[L]{\small #1\enspace\textbar\enspace #2}%
  \fancyhead[R]{\small #3: #4}%
}

% Minimalisticke cernobile boxy. Kazda definice zvlast.
\tcbset{commonlesson/.style={enhanced,breakable=false,colback=white,colframe=black,
  colbacktitle=white,coltitle=black,boxrule=0.45pt,
  arc=0.8mm,outer arc=0.8mm,left=3mm,right=3mm,top=1.6mm,bottom=1.6mm,
  toptitle=1mm,bottomtitle=0.7mm,fonttitle=\bfseries,
  attach boxed title to top left={xshift=2.5mm,yshift=-1.15mm},
  boxed title style={colback=white,colframe=white,boxrule=0pt,left=1mm,right=1mm}}}
\newtcolorbox{zadaniBox}[1][]{commonlesson,title={Zadání},#1}
\newtcolorbox{prikladBox}[1][]{commonlesson,title={Příklad},#1}
\newtcolorbox{definiceBox}[1][]{commonlesson,title={Definice},#1}

% Spolecny box pro zadani a jeho reseni, bez mezery mezi castmi.
% Pouziti:
%   \begin{prikladReseniBox}
%     Zadani ulohy.
%     \tcblower
%     \begin{align*} ... \end{align*}
%   \end{prikladReseniBox}
% Spodni cast ma jemny sedy podklad a delici linku; lze v ni sazet
% vice radku matematiky. Obe casti maji jen jeden spolecny ramecek.
\tcbset{solutionlesson/.style={
  commonlesson,
  enhanced jigsaw,
  breakable,
  colbacklower=black!3,
  segmentation style={solid,draw=black!25,line width=0.35pt},
  fontlower=\small\color{notegray},
  before lower={\noindent\textbf{Řešení:}\enspace},
  leftlower=3mm,rightlower=3mm
}}
\newtcolorbox{prikladReseniBox}[1][]{solutionlesson,title={Příklad},#1}
\newtcolorbox{zadaniReseniBox}[1][]{solutionlesson,title={Zadání},#1}

% Pomucky pro autora pripravy: sedive poznamky; \reseni pro starsi pripravy
\newcommand{\poznamka}[1]{\par\begingroup\small\color{notegray}\itshape
  \textnormal{\textbf{Poznámka pro učitele.}}\enspace #1\par\endgroup}
\newcommand{\reseni}[1]{\par\begingroup\small\color{notegray}
  \textbf{Očekávané řešení:}\enspace #1\par\endgroup}
\newcommand{\blok}[1]{\needspace{5\baselineskip}\subsection*{#1}}
\newcommand{\vyrok}[1]{\uv{#1}}

\nastavPripravu{1. ročník}{Algebraické výrazy}{56}{Smíšené úpravy algebraických výrazů}
\begin{document}
\poznamka{45 minut: 5 min výběr strategie, 12 min dva společné příklady, 18 min samostatná práce a kontrola, 5 min diskuse o chybách, 5 min shrnutí. Náročnější dvojice a rezerva dobrovolně. Opakujeme dosavadní nástroje, nezavádíme další definice.}
\blok{Nejprve zvol strategii, potom počítej}
Při různých operacích s výrazy nemá jediný univerzální \uv{první krok} smysl. Před krácením rozkládáme na součin; před sčítáním hledáme společný jmenovatel; před dělením kontrolujeme nenulovost dělitele. \textbf{Podmínky původního zápisu určíme vždy předem.}
\begin{prikladReseniBox}[breakable=false,title={Příklad 1: dvě cesty, stejný výsledek}]
Uprav $\displaystyle\left(\frac{x}{x-1}-\frac{x+1}{x}\right)(x^2-x)$.
\tcblower
Podmínky $x\ne0,1$. Nejprve společný jmenovatel:
\[\frac{x^2-(x+1)(x-1)}{x(x-1)}\cdot x(x-1)=\frac1{x(x-1)}\cdot x(x-1)=1.\]
Ještě rychlejší je roznásobit vnější činitel a krátit jednotlivé součiny; podmínky zůstávají stejné.
\end{prikladReseniBox}
\begin{prikladReseniBox}[breakable=false,title={Příklad 2: rozklad před násobením}]
Uprav $\displaystyle\frac{x^2-4}{x^2+4x+4}\cdot\frac{x+2}{x-2}$.
\tcblower
Podmínky $x\ne-2,2$. Platí
\[\frac{(x-2)(x+2)}{(x+2)^2}\cdot\frac{x+2}{x-2}=1,\quad x\ne-2,2.\]
\end{prikladReseniBox}
\blok{Samostatná práce}
\begin{zadaniReseniBox}[breakable=false,title={Procvičení: vyber vhodné pořadí úprav}]
\begin{enumerate}[label=\alph*)]
\item $\frac{(x+1)^2-(x-1)^2}{4x}$;
\item $\frac{x^2-1}{(x+1)^2}+\frac2{x+1}$;
\item $\frac{x^3-8}{x^2-4}-\frac4{x+2}$;
\item $\left(\frac1{x-1}+\frac1{x+1}\right):\frac{x}{x^2-1}$;
\item $\frac{(x-2)(x+2)}{x-2}-\frac{4}{x+2}$.
\end{enumerate}
\tcblower
\begin{enumerate}[label=\alph*)]
\item $1$, $x\ne0$; rozdíl druhých mocnin je $4x$.
\item $1$, $x\ne-1$; po krácení $\frac{x-1}{x+1}+\frac2{x+1}=1$.
\item $x$, $x\ne-2,2$; po zkrácení $x-2$ čitatel $x^2+2x+4-4=x(x+2)$.
\item $2$, $x\ne-1,0,1$; součet je $2x/(x^2-1)$.
\item $x+2-\frac4{x+2}$, $x\ne-2,2$; případně $\frac{x(x+4)}{x+2}$ při týchž podmínkách.
\end{enumerate}
\end{zadaniReseniBox}
\poznamka{Při kontrole chtěj zdůvodnění, proč zrovna tento krok zjednodušuje další výpočet. Jeden správný výsledek pro náhodné $x$ nestačí jako důkaz identity.}
\blok{Volitelná rezerva}
\begin{zadaniReseniBox}[breakable=false,title={Rozšíření: více proměnných}]
Pro $a,b\in\mathbb R$ uprav $\displaystyle\frac{a^2-b^2}{a-b}:\frac{a+b}{a}$ a určete podmínky.
\tcblower
Původní zápis vyžaduje $a\ne b$, $a\ne0$ a $a+b\ne0$. Po rozkladu $(a-b)(a+b)$ a převrácení dělitele dostaneme $a$ při všech třech podmínkách.
\end{zadaniReseniBox}
\blok{Domácí procvičení}
\begin{zadaniReseniBox}[breakable=false,title={Domácí úkol: souhrnné úpravy}]
\begin{enumerate}[label=\alph*)]
\item $\frac{x^2-9}{(x+3)^2}\cdot\frac{x+3}{x-3}$;
\item $\left(\frac{x}{x+1}-\frac{x-1}{x}\right)(x^2+x)$;
\item $\frac{x^3+8}{x^2-4}-\frac4{x-2}$;
\item $\dfrac{\frac1x-1}{\frac1x+1}$;
\item $\frac{1}{x-2}-\frac{1}{x+2}$;
\item $\left(\frac1{x-1}+\frac1{x+1}\right)\cdot(x^2-1)$.
\end{enumerate}
\tcblower
\begin{enumerate}[label=\alph*)]
\item $1$, $x\ne-3,3$.
\item $1$, $x\ne0,-1$.
\item $x$, $x\ne-2,2$.
\item $\frac{1-x}{1+x}$, $x\ne0,-1$.
\item $\frac4{x^2-4}$, $x\ne-2,2$.
\item $2x$, $x\ne-1,1$.
\end{enumerate}
\end{zadaniReseniBox}
\end{document}
